\section{Run-composition formula for \texorpdfstring{\seq{A000530}}{A000530}}
\label{sec:a000530}

For a binary word \(w\), let \(c_s(w)\) be the number of occurrences of
\(s\in\{0,1\}\), and let \(m_s(w)\) be the length of its longest run of
\(s\).  The sequence \seq{A000530} counts first-hit words ending in zero: no
proper prefix satisfies
\[
 c_0+m_0\ge 2n\qquad\text{or}\qquad c_1+m_1\ge2n,
\]
and the full word satisfies the zero condition.  Put \(N=2n-1\).  A word is
\emph{safe} when
\begin{equation}
 c_s(w)+m_s(w)\le N\qquad(s=0,1).
 \label{eq:a000530-safe}
\end{equation}

For \(r\ge1\), define
\begin{equation}
 q_N(r)=\#\left\{(x_1,\ldots,x_r)\in\Z_{>0}^r:
       \sum_{i=1}^r x_i+\max_i x_i\le N\right\},
 \label{eq:a000530-q}
\end{equation}
and set \(q_N(0)=1\).

\begin{theorem}[Run-composition formula]
\label{thm:a000530-run}
For every \(n\ge1\),
\begin{equation}
 \boxed{
 a(n)=1+\sum_{r\ge1}q_N(r)^2+
          \sum_{r\ge0}q_N(r)q_N(r+1),\qquad N=2n-1.}
 \label{eq:a000530-run}
\end{equation}
Both sums are finite; in fact \(q_N(r)=0\) for \(r\ge N\).
\end{theorem}

\begin{proof}
Safe words form the internal vertices of a finite rooted binary tree, with the
empty word at the root.  Appending either bit to a safe word either remains
safe or produces a first hit.  If the tree has \(S_n\) internal vertices, then
it has \(S_n+1\) terminal children.  Bit complementation is a fixed-point-free
bijection between terminal words whose first hit is caused by zero and those
whose first hit is caused by one.  Therefore
\begin{equation}
 a(n)=\frac{S_n+1}{2}.
 \label{eq:a000530-leaves}
\end{equation}

A symbol appearing in exactly \(r\) runs has an ordered run-length vector
\((x_1,\ldots,x_r)\).  Condition \cref{eq:a000530-safe} for that symbol is
precisely the inequality in \cref{eq:a000530-q}.  A nonempty binary word has
either the same number \(r\) of zero-runs and one-runs, with two choices of
first symbol, or \(r+1\) runs of one symbol and \(r\) of the other, with two
choices for the symbol having more runs.  Hence
\[
 S_n=1+2\sum_{r\ge1}q_N(r)^2+
        2\sum_{r\ge0}q_N(r)q_N(r+1).
\]
Substitution in \cref{eq:a000530-leaves} proves
\cref{eq:a000530-run}.

Finally, every positive \(r\)-tuple has sum at least \(r\) and maximum at
least one, so \(q_N(r)=0\) when \(r+1>N\); the slightly weaker statement
\(q_N(r)=0\) for \(r\ge N\) suffices.
\end{proof}

\subsection{Explicit finite binomial sum}

For integers \(r,m,T\), define
\begin{equation}
 B(r,m,T)=\sum_{j=0}^{r}(-1)^j\binom{r}{j}
                         \binom{T-jm}{r},
 \label{eq:a000530-B}
\end{equation}
where \(\binom{u}{r}=0\) if \(u<r\) or \(u<0\).

\begin{proposition}
\label{prop:a000530-q-explicit}
For \(r\ge1\),
\begin{equation}
 q_N(r)=
 \sum_{m=1}^{\lfloor(N-r+1)/2\rfloor}
 \left(B(r,m,N-m)-B(r,m-1,N-m)\right).
 \label{eq:a000530-q-explicit}
\end{equation}
Consequently, \cref{eq:a000530-run,eq:a000530-B,eq:a000530-q-explicit}
form an explicit finite binomial-sum formula for \(a(n)\).
\end{proposition}

\begin{proof}
Fix a possible maximum part \(m\).  The number of positive \(r\)-tuples with
all parts at most \(m\) and total sum at most \(T\) equals \(B(r,m,T)\).
Indeed, subtract one from each positive part, introduce a nonnegative slack
variable for the unused total, and apply inclusion-exclusion to the \(r\)
upper bounds.  The difference
\(B(r,m,N-m)-B(r,m-1,N-m)\) therefore counts the tuples with maximum exactly
\(m\) and total sum at most \(N-m\).  A positive \(r\)-tuple with maximum exactly \(m\) has total sum at least \(m+r-1\).  Hence it can occur only when \(2m+r-1\le N\), or \(m\le\lfloor(N-r+1)/2\rfloor\).  Summing over \(m\) gives \cref{eq:a000530-q-explicit}.
\end{proof}

\subsection{Finite-state cross-check}

The state
\[
 (c_0,c_1,m_0,m_1,\text{last symbol},\text{current run length})
\]
determines every future transition and whether a predicate first becomes true.
Equal states may therefore be merged with integer multiplicity.  The dynamic
program is finite: if both symbols occur in a safe word, then
\(c_s\le2n-2\) for each \(s\), so its length is at most \(4n-4\); constant
safe words are shorter.  Thus exhausting the safe frontier is exact rather
than cutoff-based.

The run formula and the independent state dynamic program agree through
\(n=50\), while literal word enumeration verifies the first five terms.  The
formula begins
\begin{equation*}
\begin{aligned}
1,&\ 5,\ 28,\ 226,\ 2077,\ 20770,\ 222884,\ 2529541,\\
&29995812,\ 368025335,\ 4637331008,\ 59687103424,\\
&781671686963,\ 10386703092609,\ 139745901077363,\ldots.
\end{aligned}
\end{equation*}
The stored formula sample continues through \(n=100\); the all-\(n\) formula
itself has no computational endpoint.
